\uuid{U9hS}
\titre{Équations complexes}
\auteur{bodin ; hanani ; mzari}
\organisation{exo7}
\licence{CC BY-NC-SA 4.0}
\datecreate{2020-01-20}
\niveau{L1}
\module{Algèbre}
\chapitre{Nombres complexes}
\sousChapitre{Racine carrée, équation du second degré}
\difficulte{2}
\theme{Équations, équations complexes}
\exo7qcmid{lille-1.nombres-complexes.021}

\contenu{
\begin{question}
\qclassification{104.02,104.03,104.04}
\qtags{motcle=équations complexes}
On considère l'équation : $(E) : \, z^2 = \frac{1+i}{\sqrt 2}$, $z\in \Cc$.   Quelles sont les assertions vraies ?
\begin{answers}

    \bad{Si $z$ est une solution de $(E)$, $\arg z = \frac{\pi}{8} [2\pi]$.}

    \good{Les solutions de $(E)$ sont :  $z=e^{i\frac{\pi}{8}}$ et $z=-e^{i\frac{\pi}{8}}$.}

    \good{$\cos(\frac{\pi}{8})= \frac{1}{2}\sqrt{2+\sqrt2}$ et
   $\sin(\frac{\pi}{8})= \frac{1}{2}\sqrt{2-\sqrt2}$}

    \bad{$\cos(\frac{\pi}{8})= \frac{1}{2}\sqrt{2-\sqrt2}$ et
   $\sin(\frac{\pi}{8})= \frac{1}{2}\sqrt{2+\sqrt2}$}

\end{answers}
\begin{explanations}
Utiliser l'écriture géométrique et algébrique pour résoudre l'équation et identifier la partie réelle et la partie imaginaire.
\end{explanations}
\end{question}
}
